\documentclass[12pt]{ctexart}

\usepackage{amsmath}
\usepackage{graphicx,subfigure}
\usepackage{bm}
\usepackage{geometry}
\geometry{top=2cm,left=1.5cm, right=1.5cm, bottom=2cm}
\usepackage[shortlabels]{enumitem}

\begin{document}

\title{\textbf{实验四\ 模拟示波器的使用}}
\author{1900011604 张植竣}
\date{2020年9月23日}

\maketitle

\section{实验数据和现象以及分析、处理、与结论}

\subsection{电压、时间和频率测量}

\begin{table}[ht!]
	\begin{center}
		\begin{tabular}{|c|c|c|}
			\hline
			$\boldsymbol{t_0}\mathbf{/s}$ & $\boldsymbol{K}\mathbf{/V}$ & $\boldsymbol{T}\mathbf{/ms}$ \\
			\hline
			每格扫描时间 & 偏转因数 & 从图中标尺读出的周期 \\
			\hline
			2 & 2 & 10.0 \\
			\hline
			$\boldsymbol{U_{pp}}\mathbf{/V}$ & $\boldsymbol{f}\mathbf{/Hz}$ & $\boldsymbol{U_e}\mathbf{/V}$ \\
			\hline
			从图中标尺读出的电压峰—峰值 & 由$T$算出的频率 & 由$U_{pp}$算出的电压有效值 \\
			\hline
			10.0 & 100 & 3.5 \\
			\hline
			$\boldsymbol{f'}\mathbf{/Hz}$ & $\boldsymbol{\Delta U}\mathbf{/V}$ & $\boldsymbol{\Delta t}\mathbf{/ms}$\\
			\hline
			从屏幕读出的频率 & 从读出功能算出的电压峰—峰值 & 从读出功能算出的周期 \\
			\hline
			100.00 & 10.00 & 10.06 \\
			\hline
		\end{tabular}
		\caption{周期、频率、电压峰—峰值、电压有效值等的测量数据}
	\end{center}
\end{table}

\subsubsection*{测量仪器最小分度与测量结果有效数字的关系}

\noindent 1. 测量仪器最小分度（即一“小格”）为1时，从十分位开始估读，误差位在十分位，有效数字比最小分度多一位。

\noindent 2. 测量仪器最小分度（即一“小格”）不为1时，从本位开始估读，误差位在本位，有效数字和最小分度位数相同。

\subsection{观测李萨如图形}

\begin{figure}[ht!]
	\centering
	\subfigure[$f_x:f_y=1:2$]{\includegraphics[width=0.25\linewidth]{1：2.png}}
	\subfigure[$f_x:f_y=3:2$]{\includegraphics[width=0.25\linewidth]{3：2.png}}
	\subfigure[$f_x:f_y=4:5$]{\includegraphics[width=0.25\linewidth]{4：5.png}}
	\caption{三种频率比的李萨如图形}
\end{figure}

\begin{table}[ht!]
	\begin{center}
		\begin{tabular}{|c|c|c|}
			\hline
			$\boldsymbol{f_x}\mathbf{/Hz}$ & $\boldsymbol{f_y}\mathbf{/Hz}$ & $\boldsymbol{f_x:f_y}$ \\
			\hline
			1000 & 2000 & 1:2 \\
			\hline
			171771 & 114514 & 3:2 \\
			\hline
			1000 & 1250 & 4:5 \\
			\hline
		\end{tabular}
		\caption{三种频率比的李萨如图形的相关数据}
	\end{center}
\end{table}

\subsubsection*{根据李萨如图形确定两信号的频率比}

如果两个简谐信号的频率比接近简单整数比，则可以数出水平线与图形相交的点数$m$和垂直线与图形相交的点数$n$，确定两信号频率比为:
\[
	f_y:f_x = m:n
\]
这样如果其中一个简谐信号的频率已知，就可以计算出另一信号的频率。

\newpage
\section{实验收获}

\subsection{部分示波器操作总结}

\noindent 1. 接入电路之后需要检查\ \fbox{SOURCE}\ 键是否在CH1或CH2档位，连接到这两个档位才是内触发，在LINE或EXT模式下都需要外接触发源。

\noindent 2. 调节\ \fbox{TRIG LEVEL}\ 键可以实现触发扫描同步，此时上方的触发指示灯\ \fbox{TRIG'D}\ 会亮起，如果事先知道输入信号的电压峰值，可以在屏幕上观察到触发电平值，并调节到大于半格的程度。如果触发电平高于电压峰值，会在屏幕上看到同时显示的有一定相位差的多个波形重叠成网的形状。

\noindent 3. 如果屏幕上显示了三条图线，注意\ \fbox{CH1}\ 部分的\ \fbox{ADD}\ 键是否开启。这个键可以显示CH1与CH2两信号叠加后的图像。

\noindent 4. 面板上\ \fbox{CH2}\ 部分的\ \fbox{INV}\ 键可以显示CH2信号波形反相后的图像。

\noindent 5. 使用读出功能需要开启\ \fbox{$\Delta\mathrm{V-}\Delta\mathrm{t-OFF}$}\ 键，选择测量对象为$\Delta\mathrm{V}$时为电压测量，选择测量对象为$\Delta\mathrm{t}$时为时间测量。按\ \fbox{TCK/C2}\ 键可以选择移动的光标为光标1、光标2，或跟踪（两个光标同步调节）。调节\ \fbox{TCK/C2}\ 键可以移动光标，但是光标运动不连续，可能造成较大误差。旋转一格时光标移动分度值的1/100，按一格时光标移动分度值的1/4。

\subsection{对李萨如图形观测误差的探究}

根据李萨如图形确定两信号的频率比有一定误差，经过探究发现即使两信号的频率之比不严格等于简单整数比，在示波器上显示的图像依然是规则的李萨如图形。

可以总结出以下规律：两信号频率比越接近简单整数比，李萨如图形越接近规则图形；两信号频率比越远离简单整数比，李萨如图形会趋于填满整个矩形。这两者之间存在一个“拐点”，即若$f_y:f_x = m:(n-\delta)$，形成图形的混乱程度会随着相对于简单整数比的偏离$|\delta|$的减小先增大后减小。如下图：

\begin{figure}[ht!]
	\centering
	\subfigure[$f_x:f_y=1:1.6$]{\includegraphics[width=0.2\linewidth]{1：1.6.png}}
	\subfigure[$f_x:f_y=1:1.9$]{\includegraphics[width=0.2\linewidth]{1：1.9.png}}
	\subfigure[$f_x:f_y=1:1.9999$]{\includegraphics[width=0.2\linewidth]{1：1.9999.png}}
	\subfigure[$f_x:f_y=1:2$]{\includegraphics[width=0.2\linewidth]{1：2.0.png}}
	\caption{根据李萨如图形确定两信号的频率比}
\end{figure}

即根据李萨如图形确定两信号的频率比有一定误差。实际上这个误差不仅体现在图形的样式上，也体现在李萨如图形看上去缓慢旋转的现象中。这可以说明实验中信号源、示波器等组成的电路有一定的频率不稳定性。

\end{document}